\documentclass[tikz,border=0pt]{standalone}
\usepackage[T1]{fontenc}
\usepackage{helvet}
\usepackage{amsmath,amssymb}
\usepackage{fix-cm}
\renewcommand{\familydefault}{\sfdefault}
\usetikzlibrary{arrows.meta,calc,decorations.pathreplacing,patterns,positioning}
\pdfinfoomitdate=1
\pdftrailerid{}
\definecolor{ink}{HTML}{000000}
\definecolor{blue}{HTML}{3B86FF}
\definecolor{green}{HTML}{00A14B}
\definecolor{red}{HTML}{FC4D30}
\definecolor{purple}{HTML}{A93BFF}
\definecolor{gold}{HTML}{FFC943}
\definecolor{palegold}{HTML}{FFDD92}
\definecolor{muted}{HTML}{747474}
\definecolor{neutral}{HTML}{F1F2F2}
\definecolor{rule}{HTML}{D9D9D9}
\newenvironment{candidate}[1]{%
 \begin{tikzpicture}[x=1cm,y=1cm,line width=.65pt,text=ink,
  font=\fontsize{10}{11.7}\selectfont,
  every node/.style={inner sep=1.4pt},
  flow/.style={-{Stealth[length=1.6mm,width=1.1mm]},line width=.75pt},
  softflow/.style={-{Stealth[length=1.45mm,width=1mm]},line width=.6pt},
  state/.style={circle,draw=blue,fill=blue!8,minimum size=.60cm,line width=.8pt},
  gate/.style={draw=ink,fill=neutral,rounded corners=1.8pt,line width=.65pt},
  observable/.style={draw=ink,fill=gold,rounded corners=1.2pt,line width=.65pt},
  note/.style={text=muted,font=\fontsize{10}{11.7}\selectfont}]
  \useasboundingbox (0,0) rectangle (8.60,5.60);
  \fill[white] (0,0) rectangle (8.60,5.60);
  \node[anchor=west,font=\bfseries\fontsize{12}{13}\selectfont] at (.08,5.37) {#1};
}{\end{tikzpicture}}
\newcommand{\meter}[3]{%
 \begin{scope}[shift={(#1,#2)}]
  \draw[ink,fill=#3,rounded corners=1pt] (-.24,-.22) rectangle (.24,.22);
  \draw[ink,line width=.5pt] (-.16,-.075) arc[start angle=155,end angle=25,radius=.18];
  \draw[ink,line width=.6pt,-{Stealth[length=1.0mm]}] (0,-.075)--(.105,.12);
 \end{scope}}

\begin{document}
\begin{candidate}{b}
% Universal scalar history error, not the error of a particular trained QELM.
\fill[green!8] (1.09,.97) rectangle (7.60,1.91);
\draw[flow] (1.09,.97)--(7.85,.97);
\draw[flow] (1.09,.97)--(1.09,4.93);
\node[rotate=90] at (.34,3.00) {best possible history error};
\node[anchor=north] at (4.77,.47) {window length $d$};
\draw[muted,densely dashed,line width=.6pt] (1.09,1.91)--(7.60,1.91);
\node[anchor=east] at (.96,1.91) {$\varepsilon$};
\node[text=green] at (6.27,1.61) {within tolerance};
\node[observable,minimum height=.52cm] at (5.72,4.79) {scalar optimum: $D_d/2$};
\draw[blue,line width=1.05pt] (1.31,4.57)
 .. controls (2.04,3.08) and (2.99,2.20) .. (3.94,1.91)
 .. controls (4.83,1.24) and (6.05,1.10) .. (7.50,1.07);
\draw[red,line width=1.05pt] (1.31,4.57)
 .. controls (2.17,2.59) and (3.77,2.46) .. (7.50,2.45);
\node[text=red] at (6.26,2.85) {persistent floor};
\draw[green,densely dashed,line width=.6pt] (3.94,.97)--(3.94,1.91);
\fill[green] (3.94,1.91) circle (.055);
% The marker is a conceptual crossing, not the first strict integer d_epsilon.
\node at (3.94,.73) {crossing};
\node[text=blue] at (4.60,3.83) {past influence fades};
\end{candidate}
\end{document}
